\subsection{\emph{Rucho v.\ Common Cause} and the Federal Foreclosure}

In \emph{Rucho v.\ Common Cause} (2019), the U.S. Supreme Court held 5--4
that federal courts lack jurisdiction to adjudicate partisan gerrymandering
claims under the Constitution \citep{rucho2019}.  Chief Justice Roberts'
majority opinion acknowledged that partisan gerrymandering is ``incompatible
with democratic principles'' but concluded that the ``fairness'' of electoral
maps presents a political question beyond the reach of federal judicial
standards.  The dissent (Kagan, J.) argued that manageable standards were
available, including the Efficiency Gap.

\emph{Rucho}'s practical effect is to foreclose federal constitutional
challenges to partisan gerrymandering.  State courts and state constitutions
have partially filled this gap: the North Carolina Supreme Court struck down
NC's legislative district plan on state constitutional grounds in 2022
(\emph{Harper v.\ Hall}), though subsequent court composition changes
complicated enforcement.  Pennsylvania's Supreme Court struck its
congressional map in 2018 (\emph{League of Women Voters v.\ Pennsylvania}).

\subsection{PSDs as a Legal Standard}

PSDs provide a concrete standard for measuring partisan gerrymandering that
does not require courts to evaluate ``fairness'' in the abstract---a concern
at the heart of \emph{Rucho}'s justiciability holding.  Instead, courts can
ask: does the enacted map produce more safe seats for one party than a
neutral PSD map would?  This question is empirical, not normative.

The PSD comparison has several advantages as a legal standard:

\paragraph{Algorithmic neutrality.}
PSDs do not target any specific partisan outcome.  The algorithm accepts
geographic data (partisan vote shares) and produces the most homogeneous
partition of that data.  Whichever party's voters are more geographically
sorted gets more safe seats---this is geography, not design.

\paragraph{Quantifiability.}
The EG difference between an enacted plan and a PSD baseline (e.g., $+0.14$
vs.\ $-0.01$ in NC) provides a precise measure of incremental partisan
manipulation.  Courts can establish thresholds above which the difference is
attributable to intentional manipulation rather than geographic sorting.

\paragraph{Comparability.}
Because PSDs and enacted plans are generated from the same underlying data
(same census tracts, same election returns), comparisons are apples-to-apples:
any difference in safe seat counts or EG cannot be attributed to data choices
or methodology.

\subsection{Normative Neutrality}

It is important to emphasise what PSDs are not.  PSDs are not a \emph{reform
proposal}: we do not argue that safe seats are desirable.  They are not a
\emph{gerrymandering tool}: because the algorithm distributes safe seats
proportionally between parties, PSDs would not provide a partisan advantage.
They are an \emph{analytical instrument}---a way to quantify how much of the
safe-seat concentration in an enacted plan is attributable to geography versus
intentional manipulation.

This normative neutrality is consistent with prior work that uses algorithmic
methods to evaluate redistricting plans without advocating for any specific
political outcome \citep{dellib2026edgeweighted, chen2013unintentional}.

\subsection{Limitations}

Three limitations bound this analysis.  First, all results are single-run:
ensemble variation in PSD maps is not characterised, and the reported
homogeneity levels may not be representative of the full distribution of
possible PSD maps.  Second, the 2020 presidential election results are used
as the partisan signal; different elections or multi-election averages would
produce different weight matrices and different PSD configurations.  Third,
the analysis holds geographic units fixed at census tracts; the MAUP
analysis (C.1) suggests that block-group-level results would be similar, but
this has not been tested for PSD weighting specifically.

\subsection{Relationship to Other Metrics}

The PSD framework is complementary to existing measures:

\begin{itemize}
  \item \textbf{Efficiency Gap}: PSDs produce near-zero EG by construction,
    establishing the EG a geographically sorted map would have without
    intentional manipulation.
  \item \textbf{Mean-median difference}: PSDs with near-zero EG also tend to
    produce small mean-median differences, providing a cross-validation.
  \item \textbf{Declination}: The geometric declination metric (efficiency
    of seat-vote curves) is expected to be near zero for PSD maps and can
    be compared to enacted-plan declination.
\end{itemize}
